% ==================================================================
\section{The File Drawer and Its Growing Costs}
\label{sec:background}
% ==================================================================

\subsection{The Scale of Publication Bias}

The extent of positive filtering has been quantified across multiple
disciplines. Fanelli analyzed over 4,600 papers published between
1990 and 2007 and found that the proportion reporting positive
results exceeded 80\% after 1999, peaking at 88.6\% in 2005
\citep{fanelli2012negative}. In one subfield where this pattern has
been directly quantified, McGreivy and Hakim reviewed 82 articles on
ML methods for solving fluid-related PDEs and found that 93\% claimed
superiority over traditional numerical methods, yet 79\% of those
claims relied on weak baselines. The authors concluded that negative
results were absent not because ML almost always outperforms, but
because researchers almost never publish when it does not
\citep{mcgreivy2024replacing}. Scheel et al. provided a controlled
test by comparing standard psychology papers with Registered Reports,
a format in which publication is guaranteed regardless of outcome.
The positive result rate dropped from 96\% (N=152) to 44\% (N=71)
\citep{scheel2021excess}.

\subsection{Why Failure Matters}

Research fails far more often than it succeeds. Over 90\% of drug
candidates entering clinical trials fail to reach approval
\citep{sun2022clinical}. Most hypotheses in basic science do not
survive experimental testing. The knowledge gained from these
failures is substantial. Experienced researchers accumulate failed
attempts as tacit knowledge, building an internal distribution of
prior experience that sharpens judgment about which directions are
promising and which are dead ends. This is what makes a senior
scientist valuable, and what makes doctoral training long and
difficult. Years of absorbing unpublished failure, at significant
personal cost, gradually build the intuition that no textbook can
substitute.

\subsection{The Collective Cost of Not Sharing Failure}

At the collective level, this tacit knowledge transfers informally
through mentoring and conversation, but not through the published
literature, which records only what works. Every researcher has
experienced the moment of discovering, months or years into a
project, that someone elsewhere had already attempted the same
approach and quietly abandoned it. In published literature, this
appears as being ``scooped.'' But when the prior work failed and
was never published, there is no record to be scooped by, and the
redundancy remains invisible. Chalmers and Glasziou estimated that
85\% of total research investment is avoidably wasted, with failure
to account for prior evidence, including unpublished negative
results, identified as a primary driver
\citep{chalmers2009avoidable}.

These costs have been compounding. Over the past two decades,
research output has expanded steadily, multiplying both the volume
of redundant failure and the pressure on individual researchers to
produce publishable results. At the extreme end of this pressure,
questionable research practices, from selective reporting to data
manipulation, have become more prevalent in more competitive
academic environments \citep{fanelli2010pressures}. Suspected paper
mill articles are now doubling every 1.5 years, roughly ten times
faster than the total publication base \citep{richardson2025pnas}.
In 2023 alone, over 10,000 papers were retracted, and retraction
rates have more than tripled over the past decade
\citep{noorden2023retraction}. Yet the arrangement persisted,
because recording and retrieving failure at scale remained
impractical at human reading speed, and no alternative was
economical enough to justify the shift.